\section{Objetivo dos shaders}

Os shaders apresentados nesta etapa descrevem a lógica necessária para
implementar o efeito de UV Mapping Dinâmico e Procedural em uma arquitetura
gráfica baseada em Vulkan. O pseudocódigo utiliza uma sintaxe semelhante à
linguagem GLSL, porém não tem como objetivo ser compilado diretamente.

Foram utilizados dois shaders principais: um Vertex Shader, responsável pelo
processamento dos vértices e pela preparação das coordenadas UV, e um
Fragment Shader, responsável pela aplicação do deslocamento dinâmico, pela
amostragem da textura, pelo cálculo procedural e pela produção da cor final.

O fluxo geral pode ser representado por:

\begin{center}
\begin{tabular}{c}
\textit{Vertex Buffer}
$\rightarrow$
\textit{Vertex Shader}
\\[0.2cm]
$\rightarrow$
\textit{Rasterização e Interpolação}
$\rightarrow$
\textit{Fragment Shader}
\\[0.2cm]
$\rightarrow$
\textit{Color Attachment}
$\rightarrow$
\textit{Framebuffer}
\end{tabular}
\end{center}


\section{Vertex Shader}

O Vertex Shader recebe os atributos associados aos vértices da malha. São
utilizados a posição, a normal e as coordenadas UV. O shader também recebe
por meio de um Uniform Buffer os parâmetros necessários para a transformação
geométrica e para o comportamento dinâmico das coordenadas de textura.

Os principais dados utilizados são:

\begin{itemize}
    \item \textbf{inPosition}: posição do vértice no espaço do modelo;
    \item \textbf{inNormal}: vetor normal associado ao vértice;
    \item \textbf{inUV}: coordenadas de textura do vértice;
    \item \textbf{model}: matriz de transformação do espaço do modelo para o
    espaço de mundo;
    \item \textbf{view}: matriz de transformação para o espaço de visualização;
    \item \textbf{projection}: matriz de projeção;
    \item \textbf{time}: tempo utilizado para produzir o movimento da textura;
    \item \textbf{uvVelocity}: velocidade de deslocamento das coordenadas UV;
    \item \textbf{uvScale}: escala aplicada às coordenadas UV.
\end{itemize}

\subsection{Transformação da posição}

A posição do vértice é transformada pelas matrizes Model, View e Projection.
De forma simplificada:

\begin{equation}
P_{world}=M_{model}P_{local},
\end{equation}

\begin{equation}
P_{view}=M_{view}P_{world},
\end{equation}

\begin{equation}
P_{clip}=M_{projection}P_{view}.
\end{equation}

O resultado é utilizado para determinar a posição do vértice na tela.

\subsection{Escala das coordenadas UV}

A escala permite controlar a repetição da textura na superfície. A operação
é representada por:

\begin{equation}
UV_{scale}=UV\cdot S,
\end{equation}

em que $UV$ representa a coordenada original e $S$ representa o fator de
escala.

\subsection{Deslocamento temporal}

O deslocamento dinâmico utiliza o tempo de execução e uma velocidade definida
pela aplicação:

\begin{equation}
UV_{offset}=Vt,
\end{equation}

em que $V$ representa a velocidade de deslocamento e $t$ representa o tempo.

A coordenada enviada para o próximo estágio é então:

\begin{equation}
UV'=UV_{scale}+UV_{offset}.
\end{equation}

Essa operação permite produzir movimento contínuo da textura sem modificar
a geometria da malha.

\section{Pseudocódigo do Vertex Shader}

O arquivo \texttt{uv\_mapping.vert} contém o pseudocódigo correspondente ao
Vertex Shader:

\begin{verbatim}

layout(location = 0) in vec3 inPosition;
layout(location = 1) in vec3 inNormal;
layout(location = 2) in vec2 inUV;

layout(set = 0, binding = 0) uniform Parameters
{
    mat4 model;
    mat4 view;
    mat4 projection;

    float time;
    vec2 uvVelocity;
    vec2 uvScale;
} params;

layout(location = 0) out vec2 fragUV;
layout(location = 1) out vec3 fragNormal;
layout(location = 2) out vec3 fragPosition;

void main()
{
    vec4 worldPosition =
        params.model * vec4(inPosition, 1.0);

    vec4 viewPosition =
        params.view * worldPosition;

    gl_Position =
        params.projection * viewPosition;

    fragPosition = worldPosition.xyz;
    fragNormal = inNormal;

    vec2 scaledUV =
        inUV * params.uvScale;

    vec2 uvOffset =
        params.uvVelocity * params.time;

    fragUV =
        scaledUV + uvOffset;
}
\end{verbatim}


\section{Rasterização e interpolação}

Após a execução do Vertex Shader, as primitivas são submetidas à
rasterização. Durante esse processo, os atributos enviados pelos vértices
são interpolados para os fragmentos.

Para um triângulo, a interpolação das coordenadas UV pode ser representada
por:

\begin{equation}
UV=\alpha UV_0+\beta UV_1+\gamma UV_2,
\end{equation}

com:

\begin{equation}
\alpha+\beta+\gamma=1.
\end{equation}

O resultado dessa interpolação é recebido pelo Fragment Shader como
\texttt{fragUV}.


\section{Fragment Shader}

O Fragment Shader é responsável pela etapa principal do efeito visual. Ele
recebe as coordenadas UV interpoladas, a normal e a posição do fragmento.

Também recebe os mesmos parâmetros temporais utilizados no Vertex Shader,
além dos recursos de textura.

Os principais recursos utilizados são:

\begin{itemize}
    \item \textbf{fragUV}: coordenada UV interpolada;
    \item \textbf{fragNormal}: normal interpolada;
    \item \textbf{fragPosition}: posição do fragmento;
    \item \textbf{time}: tempo da animação;
    \item \textbf{textureImage}: imagem contendo a textura;
    \item \textbf{textureSampler}: sampler utilizado para realizar a leitura;
    \item \textbf{outColor}: cor final produzida pelo shader.
\end{itemize}


\section{Amostragem da textura}

A textura é acessada utilizando as coordenadas UV resultantes. Como o efeito
dinâmico pode deslocar a coordenada para fora do intervalo $[0,1]$, o
endereçamento pode utilizar repetição.

No pseudocódigo, a operação:

\begin{equation}
UV_{sample}=fract(UV)
\end{equation}

mantém a coordenada no intervalo $[0,1)$ para representar o comportamento
de repetição.

A leitura da textura é realizada conceitualmente por:

\begin{equation}
C_{texture}=Texture(UV_{sample}).
\end{equation}

O Sampler define características da leitura, incluindo filtragem e
endereçamento. Quando mipmaps são utilizados, o nível apropriado pode ser
selecionado de acordo com a resolução aparente da textura durante a
renderização.


\section{Função procedural}

Além da textura, o Fragment Shader utiliza uma função matemática para
produzir uma variação procedural.

Neste pseudocódigo é utilizada uma função seno:

\begin{equation}
W=\sin(10U+10V+2t).
\end{equation}

O resultado da função seno pertence ao intervalo $[-1,1]$. Para convertê-lo
para o intervalo $[0,1]$, é aplicada a transformação:

\begin{equation}
W_{normalizado}=\frac{W+1}{2}.
\end{equation}

O valor resultante pode ser utilizado como intensidade para produzir uma
variação visual sobre a textura.


\section{Normalização e iluminação}

A normal recebida pelo Fragment Shader é normalizada antes de participar dos
cálculos de iluminação:

\begin{equation}
N=\frac{N}{\|N\|}.
\end{equation}

Uma direção de luz também é normalizada. O fator de iluminação é obtido pelo
produto escalar:

\begin{equation}
L=\max(0,N\cdot D),
\end{equation}

em que $N$ representa a normal normalizada e $D$ representa a direção da luz.

O uso de $\max$ evita valores negativos quando a superfície estiver voltada
para o lado oposto da fonte de iluminação.


\section{Combinação dos resultados}

O resultado procedural é convertido em uma cor:

\begin{equation}
C_{procedural}=
(W,W,W).
\end{equation}

A textura original e o resultado procedural são então combinados por meio
de uma interpolação:

\begin{equation}
C_{mix}=
(1-a)C_{texture}+aC_{procedural},
\end{equation}

em que $a$ representa o fator de mistura.

No pseudocódigo utilizado, $a=0.5$, permitindo combinar igualmente os dois
resultados.


\section{Tratamento de limites}

Antes de produzir a cor final, os componentes são limitados ao intervalo
$[0,1]$:

\begin{equation}
C_{final}=\operatorname{clamp}(C_{mix}L,0,1).
\end{equation}

Esse tratamento evita valores de cor fora do intervalo esperado.

O uso de \texttt{fract} nas coordenadas também impede que valores
excessivamente deslocados sejam enviados diretamente para a amostragem,
mantendo o comportamento de repetição da textura.


\section{Pseudocódigo do Fragment Shader}

O arquivo \texttt{uv\_mapping.frag} contém o pseudocódigo correspondente ao
Fragment Shader:

\begin{verbatim}

layout(location = 0) in vec2 fragUV;
layout(location = 1) in vec3 fragNormal;
layout(location = 2) in vec3 fragPosition;

layout(location = 0) out vec4 outColor;

layout(set = 0, binding = 0) uniform Parameters
{
    mat4 model;
    mat4 view;
    mat4 projection;

    float time;
    vec2 uvVelocity;
    vec2 uvScale;
} params;

layout(set = 0, binding = 1) uniform texture2D textureImage;
layout(set = 0, binding = 2) uniform sampler textureSampler;


float proceduralPattern(vec2 uv, float time)
{
    float wave =
        sin(
            uv.x * 10.0
            + uv.y * 10.0
            + time * 2.0
        );

    float normalizedWave =
        (wave + 1.0) * 0.5;

    return normalizedWave;
}


void main()
{
    vec2 uv = fragUV;

    vec2 sampledUV =
        fract(uv);

    vec4 textureColor =
        texture(
            sampler2D(textureImage, textureSampler),
            sampledUV
        );

    float proceduralValue =
        proceduralPattern(
            sampledUV,
            params.time
        );

    vec3 normal =
        normalize(fragNormal);

    vec3 lightDirection =
        normalize(vec3(0.4, 0.8, 0.6));

    float lightFactor =
        max(
            dot(normal, lightDirection),
            0.0
        );

    vec3 proceduralColor =
        vec3(
            proceduralValue,
            proceduralValue,
            proceduralValue
        );

    vec3 finalColor =
        mix(
            textureColor.rgb,
            proceduralColor,
            0.5
        );

    finalColor =
        finalColor * lightFactor;

    finalColor =
        clamp(
            finalColor,
            0.0,
            1.0
        );

    outColor =
        vec4(
            finalColor,
            textureColor.a
        );
}
\end{verbatim}


\section{Fluxo completo dos dados}

O processamento completo pode ser resumido da seguinte forma:

\begin{center}
\begin{tabular}{c}
\textit{Vertex Buffer}
$\rightarrow$
\textit{Vertex Shader}
\\[0.15cm]
$\downarrow$
\\[0.15cm]
\textit{Posição transformada + UV modificada}
\\[0.15cm]
$\rightarrow$
\textit{Rasterização e Interpolação}
\\[0.15cm]
$\downarrow$
\\[0.15cm]
\textit{Fragment Shader}
\\[0.15cm]
$\downarrow$
\\[0.15cm]
\textit{Sampling + Função Procedural + Iluminação}
\\[0.15cm]
$\downarrow$
\\[0.15cm]
\textit{Cor Final}
\\[0.15cm]
$\rightarrow$
\textit{Color Attachment}
\end{tabular}
\end{center}

A aplicação fornece os dados geométricos e os parâmetros dinâmicos. O
Vertex Shader transforma a geometria e prepara as coordenadas UV. A
rasterização interpola esses atributos para os fragmentos. O Fragment Shader
utiliza as coordenadas resultantes para acessar a textura e calcular a
variação procedural.

O resultado final é enviado ao Color Attachment do Render Pass, completando
o fluxo de processamento.


\section{Relação com o UV Mapping Dinâmico e Procedural}

O pseudocódigo implementa as duas características centrais estudadas neste
trabalho.

No UV Mapping Dinâmico, a coordenada de textura é alterada de acordo com o
tempo:

\begin{equation}
UV'=UV+Vt.
\end{equation}

Isso permite produzir movimento contínuo da textura sem modificar a malha.

No Mapeamento Procedural, uma função matemática utiliza a coordenada UV e o
tempo para produzir uma variação:

\begin{equation}
C=f(UV,t).
\end{equation}

A combinação dos dois mecanismos permite utilizar uma textura convencional
como base e, simultaneamente, adicionar uma variação calculada durante a
execução do shader.

Dessa maneira, a lógica descrita nos arquivos \texttt{uv\_mapping.vert} e
\texttt{uv\_mapping.frag} complementa a arquitetura Vulkan apresentada na
Entrega 2 e descreve como o efeito pode ser implementado posteriormente em
GLSL.